Dans cet exercice, les quatre questions sont indépendantes.

Il faut traiter les quatre questions.

\subsection*{Question 1}
	Soit $f$ la fonction définie sur $]-1 ;+\infty[$ par \quad $f(x)=3+5 \ln (x+1)$

	Calculer $f(0)$.

\subsection*{Question 2}
	Justifier, en détaillant les calculs, l'égalité suivante :
	\[\ln (9)+\ln (4)=2 \ln (6)\]

\subsection*{Question 3}
	Résoudre dans $\R$ l'équation :
	\[\e^{2 x}=3\]
	On donnera la valeur exacte.

\section*{Question 4}
	Un objet se déplace à la vitesse $v$, exprimée en $\mathrm{m} \cdot \mathrm{s}^{-1}$.

	À l'instant $t=\np[s]{0}$, sa vitesse initiale est $v(0)=\np[m\cdot s^{-1}]{5}$.

	Son accélération $a$ exprimée en $\mathrm{m} \cdot \mathrm{s}^{-2}$ est modélisée en fonction du temps $t$, exprimé en secondes, sur l'intervalle $[0 ; 10]$ par
	\[a(t)=\np[m\cdot s^{-2}]{2}\]
	On rappelle que $v$ est la primitive de la fonction $a$ qui vérifie la condition initiale $v(0)=5$.

	Déterminer l'expression de $v$.
